% The abstract, shared by the preprint and the journal submission.
% Neural Networks requires at most 250 words; scripts/check_abstract.py enforces it.
Papers that reverse-engineer a neural network end with a number: \emph{our
circuit recovers 87\% of the model's performance}. That number is usually the
winner of a contest. The investigator scored many candidate circuits and reported
the best one, using the same interventions to pick the winner and to score it.
Winners of contests look better than they are, and the error bar printed beside
the number is computed as though no contest had taken place. In our experiments
that error bar, nominally $95\%$ correct, is right \MyNaiveCovHundred{} of the
time when a hundred equally good circuits are compared.

We make the number trustworthy. Treating faithfulness as a population quantity
and the search as a selection rule, we derive lower confidence bounds that stay
valid however the circuit was found, and show that the inflation of a selected
score grows like the square root of the log of the number of circuits tried,
divided by the number of interventions. We give five corrections and say when
each applies, then measure coverage instead of assuming it: on five small
transformers whose algorithms we control, the true faithfulness of every
hypothesis is known exactly. On the
indirect-object-identification circuit of GPT-2 small, searching \MyIOIm{}
hypotheses with $n=1000$ interventions, the conventional interval is right
\MyIOINaiveCov{} of the time; our band is right \MyIOIBootCov{} and still
certifies \MyIOIBootLCB{}, against an inflated estimate of \MyIOINaivePoint{}.
Honesty is cheap. Two warnings follow: prompt templates, not prompts, are the
unit of replication, and adaptive searches must hold interventions out.
